\section{Conclusion}

Aiyagari's log economy has a unique stationary equilibrium, and the statement is now a theorem
about his model as he wrote it, with his parameters, for any of his earnings processes, checked
by a machine. The proof is the one the picture suggested, supply rises and demand falls, but the
constants in the published versions of that argument do not survive $\beta=0.96$, and what does
survive are consequences of the Euler inequality at the borrowing constraint: a corner condition, a
minimal marginal propensity to consume, and a Doeblin atom. None of it uses a derivative of the
value function.

For $\mu\in\{3,5\}$ we have reduced uniqueness to an inequality between two consumption functions
that holds numerically at every wealth level with an explicit margin and that no available bound
proves; for Huggett, to a quantitative race between creditors and debtors whose creditor side we
cannot bound below. We regard both as genuinely open and have tried to say precisely what a proof
would need: comparative statics of the consumption function in the interest rate, uniform on the
support in relative terms. That is a question about the income fluctuation problem, not about
equilibrium, and it may be the right next target.

On method, the experience of this development is that a proof assistant changes what one writes
down. Arguments that quantify over ``a sufficiently large cap'' or ``a small enough interval of
rates'' had to be given their constants, and at Aiyagari's parameters several of the published
constants were found to be too small by orders of magnitude; the replacements are simpler than the
originals. Differentiability, which enters every textbook treatment through the envelope theorem,
was never needed. And the exact location of the obstruction at $\mu>1$, the wasted interest-income
term in a four-line argument, was found by reading the formal proof rather than the informal one.
